\documentclass[aspectratio=169]{beamer}
\usetheme{Madrid}
\usecolortheme{default}
\usepackage[utf8]{inputenc}
\usepackage{amsmath}
\usepackage{amsfonts}
\usepackage{amssymb}

\setbeamertemplate{navigation symbols}{}
\setbeamerfont{block body}{size=\small}
\setbeamerfont{alertblock body}{size=\small}

\title{CRISC Domain B -- 100 Exam-Style Questions}
\subtitle{IT Risk Assessment -- Realistic Scenario Practice (Answer on the Next Slide)}
\author{CRISC Mastery Series}
\date{\today}

% ------------------------------------------------------------
% \hq{number}{topic}{question}{options}{letter}{explanation}{trap}
% Creates one question frame and one explanation frame.
% (No blank lines are allowed inside the arguments.)
% ------------------------------------------------------------
\newcommand{\hq}[7]{%
	\begin{frame}{Q#1: #2}
		\begin{block}{Question}
			#3
		\end{block}
		\begin{enumerate}[(A)]
			#4
		\end{enumerate}
	\end{frame}
	\begin{frame}{Q#1: Explanation}
		\begin{block}{Correct Answer: #5}
			\begin{itemize}
				#6
			\end{itemize}
		\end{block}
		\begin{alertblock}{Critical Trap}
			#7
		\end{alertblock}
	\end{frame}
}

\begin{document}
	
	\begin{frame}
		\titlepage
	\end{frame}
	
	\begin{frame}{Instructions}
		\begin{alertblock}{How to Use This Deck}
			Think like a risk practitioner: business objectives first, owners decide, assess before acting.\\
			Read for BEST, FIRST and MOST. Answer on the slide after each question.
		\end{alertblock}
	\end{frame}
	
	\begin{frame}{Outline}
		\tableofcontents
	\end{frame}
	
	% ============================================================
	\section{Risk Identification}
	% ============================================================
	
	\hq{1}{Cloud Migration}
	{A company plans to migrate its core applications to a public cloud. What should the risk practitioner do FIRST?}
	{\item Select cloud security controls
		\item Review the provider\textquotesingle s SOC report
		\item Understand the business objectives and scope of the migration
		\item Update the disaster recovery plan}
	{C}
	{\item Risk only has meaning relative to business objectives
		\item Scope tells you which assets and processes are affected
		\item Controls and assurance reviews come after the risk context is known
		\item Starting with controls skips the purpose of the change}
	{Start with objectives and scope, not controls.}
	
	\hq{2}{Emerging Risks}
	{Which is the BEST source for identifying EMERGING IT risks?}
	{\item Prior-year audit reports
		\item Industry threat intelligence and environmental scanning
		\item The existing risk register
		\item Incident logs from the last quarter}
	{B}
	{\item Emerging risks have little or no history
		\item Forward-looking sources show trends before they become incidents
		\item Audit reports, registers and incident logs show known past issues}
	{Historical sources cannot identify emerging risks.}
	
	\hq{3}{Scenario Development}
	{Risk scenarios are MOST useful when they are developed:}
	{\item With input from both business and IT stakeholders
		\item From vulnerability scan results only
		\item By the IT department alone
		\item By external consultants without staff input}
	{A}
	{\item Business owners understand impact
		\item IT understands threats and vulnerabilities
		\item Combined input produces realistic, relevant scenarios
		\item One-sided input leaves blind spots}
	{Scenarios need both business and technical views.}
	
	\hq{4}{Risk Versus Vulnerability}
	{Which of the following BEST describes a RISK (rather than a threat or vulnerability)?}
	{\item A criminal group targeting retailers
		\item A weak password policy
		\item An unpatched server
		\item Customer data could be stolen through an unpatched server, resulting in regulatory fines}
	{D}
	{\item A risk links a cause, an event and a business consequence
		\item An unpatched server and a weak policy are vulnerabilities
		\item A criminal group is a threat source}
	{Risk is cause plus event plus consequence.}
	
	\hq{5}{New Regulation}
	{A new data protection regulation is announced. What should the risk practitioner do FIRST?}
	{\item Notify all customers
		\item Rewrite all policies
		\item Determine whether and how the regulation applies to the organization
		\item Implement new controls immediately}
	{C}
	{\item Applicability must be confirmed before any response
		\item Scope determines the gap analysis
		\item Controls and policy changes follow the assessment}
	{Assess applicability before acting.}
	
	\hq{6}{Acquisition}
	{An organization acquires a company and plans to connect the two networks next week. What is the BEST action for the risk practitioner?}
	{\item Assess the risks of the acquired environment before connecting
		\item Rely on the seller	extquotesingle s assurances
		\item Wait for the next annual risk assessment
		\item Connect and monitor traffic afterward}
	{A}
	{\item Connecting networks can import unknown risks
		\item Assessment should happen before the connection
		\item Monitoring after the fact is reactive
		\item Seller assurances are not independent evidence}
	{Assess before you connect.}
	
	\hq{7}{Vulnerability Scan Results}
	{A scan reports 3,000 vulnerabilities. What should the risk practitioner do NEXT?}
	{\item Rank them by age
		\item Patch everything immediately
		\item Report the total count to the board
		\item Relate the findings to critical assets, threats and business impact}
	{D}
	{\item Raw counts do not measure risk
		\item Business impact and threat likelihood decide priority
		\item Blanket patching is rarely feasible or necessary}
	{Findings become risks only when tied to business impact.}
	
	\hq{8}{Delphi Technique}
	{Which technique BEST obtains unbiased expert consensus when strong personalities may dominate meetings?}
	{\item Walkthrough meeting
		\item Checklist review
		\item Open brainstorming
		\item Delphi technique}
	{D}
	{\item Experts respond anonymously in rounds
		\item Feedback is shared without attribution
		\item This reduces dominance and groupthink}
	{Anonymity reduces dominant-voice bias.}
	
	\hq{9}{Shadow IT}
	{Business units are using unapproved cloud tools. What should the risk practitioner do FIRST?}
	{\item Block all unapproved tools immediately
		\item Record it as accepted
		\item Determine what data and processes are involved and assess the risk
		\item Discipline the users}
	{C}
	{\item Understand exposure before reacting
		\item Blocking may disrupt legitimate business needs
		\item Acceptance requires an assessment and an authorized owner}
	{Assess before blocking or accepting.}
	
	\hq{10}{Incident Lessons}
	{After a major incident, how should risk identification be updated?}
	{\item Remove related risks because they materialized
		\item Only update the incident log
		\item No change: the incident is closed
		\item Use root-cause analysis to revise scenarios and the risk register}
	{D}
	{\item Incidents reveal gaps in earlier assessments
		\item Root cause shows which scenarios or controls were wrong
		\item The register should reflect what was learned}
	{Incidents are inputs to risk identification.}
	
	\hq{11}{Incomplete Asset Inventory}
	{What is the MOST significant effect of an incomplete asset inventory on risk assessment?}
	{\item Reports look less professional
		\item Risks to unknown assets are not identified or assessed
		\item Control testing costs more
		\item Assessments take longer}
	{B}
	{\item You cannot assess what you do not know exists
		\item Unknown assets may be critical
		\item Exposure is understated}
	{Unknown assets mean unassessed risk.}
	
	\hq{12}{Duplicate Risks}
	{Two departments logged the same risk with different ratings. What is the BEST action?}
	{\item Keep both entries
		\item Use the higher rating
		\item Remove the lower-rated entry
		\item Consolidate, reassess using common criteria and assign one owner}
	{D}
	{\item Duplicates cause confusion and conflicting actions
		\item Common criteria give a consistent rating
		\item A single owner provides accountability}
	{One risk, one owner, one rating.}
	
	\hq{13}{Top-Down and Bottom-Up}
	{Using both top-down and bottom-up risk identification BEST helps to:}
	{\item Reduce the number of risks
		\item Guarantee completeness
		\item Link strategic risks with operational detail
		\item Remove the need for workshops}
	{C}
	{\item Top-down starts from objectives
		\item Bottom-up starts from processes and assets
		\item Together they connect strategy to operations
		\item No method guarantees completeness}
	{Combined approaches improve coverage, not completeness.}
	
	\hq{14}{Scenario Relevance}
	{Which element MOST helps decide whether a risk scenario matters to the business?}
	{\item The seniority of the analyst
		\item Its potential effect on business objectives
		\item The technology involved
		\item The number of systems affected}
	{B}
	{\item Risk is about effect on objectives
		\item Technology details do not show business relevance
		\item Objective impact drives prioritization}
	{No objective at stake means no business risk.}
	
	\hq{15}{Project Initiation}
	{When should risk identification for a new IT project begin?}
	{\item After the first incident
		\item At project initiation and continue throughout
		\item After testing is complete
		\item At go-live}
	{B}
	{\item Early identification allows cheaper treatment
		\item Risks change as the project evolves
		\item Late identification limits options}
	{Identify early, then keep going.}
	
	% ============================================================
	\section{Risk Analysis}
	% ============================================================
	
	\hq{16}{Qualitative Approach}
	{Reliable historical loss data is NOT available. Which approach is MOST appropriate?}
	{\item Using a competitor	extquotesingle s data as exact values
		\item A quantitative assessment using assumed numbers, undisclosed
		\item Skipping the assessment
		\item A qualitative assessment that clearly states its limitations}
	{D}
	{\item Qualitative methods suit limited data
		\item Limitations must be disclosed
		\item Undisclosed assumptions give false precision}
	{Do not fake precision.}
	
	\hq{17}{Risk Formula}
	{Risk is BEST expressed as a function of:}
	{\item Vulnerability and control
		\item Threat and asset
		\item Cost and benefit
		\item Likelihood and impact}
	{D}
	{\item Likelihood is how probable an event is
		\item Impact is the consequence if it occurs
		\item Together they express risk magnitude}
	{Risk equals likelihood and impact.}
	
	\hq{18}{Inherent Risk}
	{Inherent risk is the risk:}
	{\item Before considering existing controls
		\item After treatment is complete
		\item Accepted by management
		\item After controls are applied}
	{A}
	{\item It shows exposure without controls
		\item It reveals how much controls are carrying
		\item Residual risk comes after controls}
	{Inherent means before controls.}
	
	\hq{19}{Residual Risk}
	{Residual risk is BEST described as:}
	{\item Risk transferred to an insurer
		\item Total risk without controls
		\item Risk that has been eliminated
		\item The risk remaining after controls and responses are applied}
	{D}
	{\item It is what remains after treatment
		\item It is compared with appetite
		\item It is never exactly zero}
	{Residual is after controls.}
	
	\hq{20}{Impact Driver}
	{Which factor MOST influences the impact rating of an IT asset failure?}
	{\item Criticality of the business processes it supports
		\item Purchase price of the asset
		\item Age of the hardware
		\item Vendor name}
	{A}
	{\item Impact is measured by business consequence
		\item Critical processes mean higher impact
		\item Cost and age are weak proxies}
	{Business criticality drives impact.}
	
	\hq{21}{Outdated Data}
	{Loss data is 5 years old and the business has changed significantly. What should the risk practitioner do?}
	{\item Use the data unchanged
		\item Use only industry averages
		\item Discard the analysis entirely
		\item Adjust the analysis using current threat information and expert judgment}
	{D}
	{\item Old data may not reflect current risk
		\item Adjust for current conditions
		\item Data is still a useful starting point}
	{Old data is not automatically relevant data.}
	
	\hq{22}{Inconsistent Ratings}
	{Different teams rate similar risks inconsistently. What is the BEST solution?}
	{\item Let each team choose its own scales
		\item Use only the highest ratings
		\item Reduce the number of teams
		\item A standardized methodology with defined likelihood and impact scales}
	{D}
	{\item Common definitions give comparable results
		\item Scales need clear descriptions
		\item Calibration workshops help consistency}
	{Standardize scales to get comparable ratings.}
	
	\hq{23}{Sensitivity Analysis}
	{Which technique BEST shows how changes in assumptions affect risk results?}
	{\item Sensitivity analysis
		\item Walkthrough
		\item Control self-assessment
		\item Gap analysis}
	{A}
	{\item It varies inputs to see the effect on results
		\item It identifies assumptions that matter most
		\item It helps judge reliability of an analysis}
	{Test assumptions with sensitivity analysis.}
	
	\hq{24}{Annual Loss Expectancy}
	{An asset is worth \$500K, the exposure factor is 40\%, and the event is expected once every 2 years. What is the annualized loss expectancy?}
	{\item \$250K
		\item \$100K
		\item \$50K
		\item \$200K}
	{B}
	{\item SLE $= 500K \times 40\% = \$200K$
		\item ARO $= 0.5$ per year
		\item ALE $= 200K \times 0.5 = \$100K$
		\item \$200K is the single loss expectancy}
	{ALE equals SLE times ARO.}
	
	\hq{25}{Low Likelihood High Impact}
	{A risk has very low likelihood but catastrophic impact. What is the BEST view?}
	{\item Do not dismiss it; consider contingency planning or transfer
		\item Avoid it by shutting down operations
		\item Accept it without review
		\item Ignore it because likelihood is low}
	{A}
	{\item Severe impact can threaten survival
		\item Contingency and insurance suit rare severe events
		\item Low likelihood is not zero}
	{Rare does not mean unimportant.}
	
	\hq{26}{Input Quality}
	{The reliability of a risk analysis depends MOST on:}
	{\item The number of analysts
		\item The complexity of the model
		\item The length of the report
		\item The quality of the data and assumptions used}
	{D}
	{\item Poor input produces poor output
		\item Assumptions should be documented and validated
		\item Complex models do not fix weak data}
	{Garbage in, garbage out.}
	
	\hq{27}{Aggregated View}
	{Why is aggregating related risks useful?}
	{\item It lowers total risk
		\item It shows combined exposure that individual ratings may hide
		\item It removes the need for risk owners
		\item It speeds up control testing}
	{B}
	{\item Related risks can share causes
		\item Combined exposure can exceed appetite even if each item is small
		\item Aggregation supports better decisions}
	{Many small related risks can add up.}
	
	\hq{28}{Equal Scores}
	{Two risks have the same score. What is the BEST way to prioritize them?}
	{\item Treat them as equal
		\item Choose alphabetically
		\item Consider business impact on critical objectives and how quickly each could materialize
		\item Fix the cheaper one first}
	{C}
	{\item Same score can hide different profiles
		\item Impact on key objectives matters
		\item Speed of onset affects response time}
	{Same score does not mean same priority.}
	
	\hq{29}{Control Effectiveness}
	{Control effectiveness is MOST directly used to determine:}
	{\item Risk appetite
		\item Asset value
		\item Inherent risk
		\item Residual risk}
	{D}
	{\item Effective controls reduce inherent risk
		\item The result is residual risk
		\item Appetite is set by leadership}
	{Controls move you from inherent to residual.}
	
	\hq{30}{Complacency}
	{A team rates a risk low because nothing has happened for years. This reflects:}
	{\item Excessive control
		\item A valid conclusion
		\item Appropriate conservatism
		\item Complacency from past experience; current threats should be considered}
	{D}
	{\item Absence of incidents is not evidence of low risk
		\item Threats and business change over time
		\item Use current intelligence}
	{Nothing happened is not a risk rating.}
	
	% ============================================================
	\section{Risk Evaluation and Appetite}
	% ============================================================
	
	\hq{31}{Setting Appetite}
	{Who is ultimately responsible for setting risk appetite?}
	{\item The risk practitioner
		\item Internal audit
		\item The board and senior management
		\item The IT manager}
	{C}
	{\item Appetite reflects strategy and risk-taking intent
		\item Leadership is accountable for it
		\item Others advise and implement}
	{Leadership sets appetite.}
	
	\hq{32}{Above Appetite}
	{A risk exceeds appetite. What should the risk practitioner do FIRST?}
	{\item Lower the rating
		\item Inform the risk owner and escalate according to governance
		\item Implement a control immediately
		\item Accept the risk}
	{B}
	{\item The owner must decide the response
		\item Escalation follows defined thresholds
		\item The practitioner advises and does not decide}
	{Escalate; do not decide for the owner.}
	
	\hq{33}{Tolerance}
	{Risk tolerance is BEST described as:}
	{\item The control budget
		\item Total risk faced
		\item The maximum loss the organization can survive
		\item The acceptable deviation from risk appetite}
	{D}
	{\item Appetite is the desired level
		\item Tolerance is the allowed variation around it
		\item Capacity is the survival limit}
	{Tolerance is the permitted variance.}
	
	\hq{34}{Residual Above Appetite}
	{After treatment, residual risk still exceeds appetite. What is the BEST next step?}
	{\item Close the risk
		\item Escalate for further treatment or formal acceptance by the appropriate authority
		\item Lower the rating
		\item Ignore it until the next review}
	{B}
	{\item Residual above appetite needs a decision
		\item Options are more treatment or authorized acceptance
		\item Quiet closure hides exposure}
	{Above appetite needs an explicit decision.}
	
	\hq{35}{Acceptance Authority}
	{Who should accept a residual IT risk?}
	{\item The vendor
		\item The CISO by default
		\item The risk practitioner
		\item The business owner with authority over the affected process}
	{D}
	{\item The owner bears the consequences
		\item Authority must match the size of the risk
		\item The practitioner informs the decision}
	{Business owners accept business risk.}
	
	\hq{36}{Measurable Appetite}
	{The appetite statement is general and has no measures. What is the BEST improvement?}
	{\item Remove it
		\item Let each business unit invent its own
		\item Define measurable thresholds and link them to KRIs
		\item Make the statement longer}
	{C}
	{\item Measures make breaches visible
		\item KRIs track exposure against thresholds
		\item Vague statements cannot guide decisions}
	{If it cannot be measured, it cannot be managed.}
	
	\hq{37}{Risk Capacity}
	{Risk capacity is:}
	{\item The desired level of risk
		\item The budget for controls
		\item The maximum risk the organization can bear without threatening survival or objectives
		\item The acceptable variation from appetite}
	{C}
	{\item Capacity is an upper limit
		\item Appetite and tolerance sit below it
		\item Exceeding it is existential}
	{Capacity is the ceiling.}
	
	\hq{38}{Project Approval}
	{How should risk be considered when approving a new project?}
	{\item Compare expected risk with risk appetite
		\item Approve if no incidents occurred last year
		\item Approve if IT agrees
		\item Approve if cost is within budget}
	{A}
	{\item Decisions should be risk-informed
		\item Appetite provides the reference point
		\item Budget and IT agreement do not show risk acceptability}
	{Use appetite as the yardstick.}
	
	\hq{39}{Rating Disagreement}
	{IT and the business disagree about an impact rating. What is the BEST approach?}
	{\item IT decides
		\item Escalate every disagreement to the board
		\item Use the average of both
		\item Use agreed criteria, with the business owner validating impact}
	{D}
	{\item The business understands impact
		\item Agreed criteria reduce subjectivity
		\item Escalation is for unresolved or material issues}
	{Impact belongs to the business.}
	
	\hq{40}{Accepted Risk}
	{A low risk has been accepted. What should happen afterward?}
	{\item Transfer it
		\item Remove it from the register
		\item Document it and monitor for changes
		\item Forget it}
	{C}
	{\item Acceptance is a documented decision
		\item Conditions can change
		\item Monitoring confirms it stays acceptable}
	{Accepted does not mean forgotten.}
	
	% ============================================================
	\section{Risk Response}
	% ============================================================
	
	\hq{41}{Insurance}
	{Buying cyber insurance is an example of risk:}
	{\item Transfer
		\item Mitigation
		\item Acceptance
		\item Avoidance}
	{A}
	{\item Insurance shifts financial consequence
		\item It does not remove the risk
		\item Accountability remains with the organization}
	{Insurance shifts financial loss only.}
	
	\hq{42}{Cost Exceeds Benefit}
	{A control would cost \$100K per year to protect against an ALE of \$40K. What is the BEST response?}
	{\item Avoid the process entirely
		\item Accept the risk after documentation, or look for a cheaper control
		\item Implement it anyway
		\item Ignore the risk silently}
	{B}
	{\item Cost exceeds benefit
		\item Acceptance must be documented and authorized
		\item A cheaper control may exist}
	{Do not spend more than the risk is worth.}
	
	\hq{43}{Outsourcing}
	{When processing is outsourced, accountability for the associated risk:}
	{\item Remains with the organization
		\item Moves to the regulator
		\item Moves to the vendor
		\item Is shared equally by contract}
	{A}
	{\item Responsibility for tasks can be delegated
		\item Accountability cannot be outsourced
		\item Oversight remains necessary}
	{You can outsource work, not accountability.}
	
	\hq{44}{Response Selection}
	{Risk response selection should BEST be driven by:}
	{\item Cost-benefit and alignment with risk appetite
		\item The cheapest option
		\item Technology preference
		\item What peers do}
	{A}
	{\item Responses must bring risk within appetite
		\item Cost must be justified by benefit
		\item Peers	extquotesingle{} practices may differ in context}
	{Fit the response to appetite and cost.}
	
	\hq{45}{Avoidance}
	{Which is an example of risk avoidance?}
	{\item Documenting acceptance
		\item Adding a firewall
		\item Discontinuing a service that creates unacceptable risk
		\item Buying insurance}
	{C}
	{\item Avoidance removes the activity
		\item Insurance is transfer
		\item A firewall is mitigation}
	{Avoidance stops the activity.}
	
	\hq{46}{Preventive Control}
	{Which is a PREVENTIVE control?}
	{\item Restoring from backup
		\item Access control that blocks unauthorized logins
		\item Incident response procedure
		\item Review of system logs}
	{B}
	{\item Preventive controls stop events
		\item Log review is detective
		\item Backup restore is corrective}
	{Preventive stops it before it happens.}
	
	\hq{47}{Detective Control}
	{Which is a DETECTIVE control?}
	{\item Data encryption
		\item Firewall rules
		\item Segregation of duties
		\item Review of system logs for anomalies}
	{D}
	{\item Detective controls identify events after they occur
		\item Firewalls and encryption are preventive
		\item Segregation of duties is preventive}
	{Detective finds, it does not stop.}
	
	\hq{48}{Compensating Control}
	{Segregation of duties is impractical in a small team. Which is the BEST compensating control?}
	{\item Disabling logging
		\item Sharing passwords to speed work
		\item Annual rotation only
		\item Independent management review of transactions}
	{D}
	{\item Compensating controls reduce the same risk another way
		\item Independent review detects misuse
		\item It must be as effective as the original control}
	{Compensate for the same risk, not just any risk.}
	
	\hq{49}{Response Approval}
	{Who should approve a risk response plan?}
	{\item The risk practitioner
		\item The external auditor
		\item The help desk
		\item The risk owner}
	{D}
	{\item The owner is accountable for the risk
		\item The practitioner advises
		\item Auditors provide independent assurance only}
	{The owner approves the response.}
	
	\hq{50}{After Implementation}
	{After a new control is implemented, what should the risk practitioner do?}
	{\item Stop monitoring
		\item Close all related risks
		\item Assume the risk is zero
		\item Validate its effectiveness and reassess residual risk}
	{D}
	{\item Implementation does not prove effectiveness
		\item Residual risk must be re-evaluated
		\item Monitoring continues}
	{Implemented does not mean effective.}
	
	\hq{51}{Verbal Acceptance}
	{A manager accepted a risk verbally. What should the risk practitioner do?}
	{\item Implement a control without approval
		\item Treat verbal acceptance as sufficient
		\item Ensure formal, documented acceptance according to policy
		\item Delete the risk}
	{C}
	{\item Acceptance must be traceable
		\item Documentation shows informed decision-making
		\item Policy defines authority and format}
	{Undocumented acceptance is not defensible.}
	
	\hq{52}{Vendor Refuses Assurance}
	{A critical vendor refuses to provide assurance reports. What is the BEST action?}
	{\item Escalate to the risk owner and seek alternative assurance or compensating controls
		\item Accept the risk automatically
		\item Proceed without review
		\item Terminate the vendor immediately}
	{A}
	{\item Lack of assurance means unknown risk
		\item The owner decides how to respond
		\item Alternatives may include contractual rights or added controls}
	{No assurance means escalate, not ignore.}
	
	\hq{53}{Temporary Exception}
	{A temporary exception to policy is requested. Which condition is BEST?}
	{\item Open-ended with no review
		\item Time-bound, formally approved and supported by compensating controls
		\item Verbally approved
		\item Approved without compensating controls}
	{B}
	{\item Exceptions should expire
		\item Approval creates accountability
		\item Compensating controls limit exposure}
	{Exceptions must be temporary and controlled.}
	
	\hq{54}{Response Underperforms}
	{A risk response is not reducing risk as expected. What should be done?}
	{\item Stop monitoring the risk
		\item Continue unchanged
		\item Reassess the risk and revise the response
		\item Accept the risk without analysis}
	{C}
	{\item The assumed benefit was not achieved
		\item Reassessment shows true residual risk
		\item Revision aligns response with appetite}
	{A failing response needs rework.}
	
	\hq{55}{Mitigation Fit}
	{A risk exceeds appetite and an affordable, effective control exists. Which response BEST fits?}
	{\item Mitigate by implementing the control
		\item Accept the risk
		\item Transfer only, with no other action
		\item Ignore the risk}
	{A}
	{\item Risk is above appetite
		\item A cost-effective control is available
		\item Mitigation brings the risk within appetite}
	{Use available effective controls to meet appetite.}
	
	% ============================================================
	\section{Controls}
	% ============================================================
	
	\hq{56}{Design Versus Operation}
	{A control is well designed but staff do not perform it. This is a deficiency in:}
	{\item Design
		\item Operating effectiveness
		\item Risk appetite
		\item Scope}
	{B}
	{\item Design is correct
		\item Failure is in execution
		\item Remediation focuses on training or process}
	{Good design does not guarantee operation.}
	
	\hq{57}{Defense in Depth}
	{Defense in depth BEST means:}
	{\item One very strong control
		\item More policies
		\item Perimeter controls only
		\item Multiple layered controls so one failure does not expose assets}
	{D}
	{\item Layers provide redundancy
		\item No single failure exposes the asset
		\item Layers should not share the same weakness}
	{Layers beat a single point of defense.}
	
	\hq{58}{Control Owner}
	{The control owner is responsible for:}
	{\item Independently auditing the control
		\item Accepting residual risk
		\item Setting risk appetite
		\item Implementing and operating the control}
	{D}
	{\item Control owners run the control
		\item Risk owners accept risk
		\item Audit is independent}
	{Control owner operates; risk owner accepts.}
	
	\hq{59}{Corrective Control}
	{Restoring systems from backup after an incident is primarily a:}
	{\item Deterrent control
		\item Corrective control
		\item Detective control
		\item Preventive control}
	{B}
	{\item Corrective controls limit damage and restore
		\item Preventive controls stop events
		\item Detective controls identify events}
	{Restoring is corrective.}
	
	\hq{60}{Independent Testing}
	{Which testing provides the MOST objective assurance of control effectiveness?}
	{\item Verbal confirmation
		\item Manager attestation without evidence
		\item A self-test by the control owner
		\item Independent testing by a party not responsible for the control}
	{D}
	{\item Independence reduces bias
		\item Evidence supports conclusions
		\item Self-assessment is useful but less objective}
	{Independence gives objectivity.}
	
	\hq{61}{Control Gap}
	{A control gap is discovered. What should the risk practitioner do FIRST?}
	{\item Report to regulators
		\item Implement a control immediately
		\item Assess the risk exposure resulting from the gap
		\item Close the gap by changing policy}
	{C}
	{\item Impact of the gap determines urgency
		\item Response should be proportionate
		\item Assessment comes before remediation}
	{Assess the exposure, then respond.}
	
	\hq{62}{Automated Controls}
	{Compared with manual controls, automated controls are generally:}
	{\item Immune to failure
		\item Not requiring testing
		\item More consistent, with less risk of human error
		\item Always cheaper}
	{C}
	{\item Automation applies rules consistently
		\item They can still be misconfigured
		\item They need testing and change control}
	{Automated is consistent, not infallible.}
	
	\hq{63}{Control Selection}
	{Controls should be selected BEST based on:}
	{\item Risk assessment results and risk appetite
		\item Vendor recommendations
		\item What peers use
		\item Cost alone}
	{A}
	{\item Controls address identified risks
		\item Appetite defines how much reduction is needed
		\item Cost is a factor but not the driver}
	{Controls follow risk, not fashion.}
	
	\hq{64}{Overlapping Controls}
	{Two controls address the same risk. This is:}
	{\item Always wasteful
		\item Acceptable if it supports defense in depth and is justified
		\item Prohibited
		\item A policy violation}
	{B}
	{\item Overlap can provide redundancy
		\item Justification should be documented
		\item Review cost and benefit}
	{Redundancy can be intentional.}
	
	\hq{65}{CSA Limitation}
	{The MAIN limitation of control self-assessment is:}
	{\item Possible bias from those assessing their own controls
		\item It is too slow
		\item It is not relevant to risk
		\item It is too costly}
	{A}
	{\item Owners may rate themselves favorably
		\item Independent validation adds assurance
		\item CSA still supports engagement}
	{Self-assessment needs validation.}
	
	% ============================================================
	\section{Monitoring and Reporting}
	% ============================================================
	
	\hq{66}{KRI Quality}
	{Which is the MOST important characteristic of a good KRI?}
	{\item It predicts changes in risk exposure and can be measured
		\item It is easy to collect
		\item It is inexpensive
		\item It never changes}
	{A}
	{\item KRIs should be forward-looking
		\item Measurability supports thresholds
		\item Ease and cost are secondary}
	{A KRI should be predictive and measurable.}
	
	\hq{67}{KRI Breach}
	{A KRI exceeds its threshold. What should happen FIRST?}
	{\item Validate the data, determine the cause and escalate per procedure
		\item Wait for the next cycle
		\item Raise the threshold
		\item Buy new controls immediately}
	{A}
	{\item Confirm the data is accurate
		\item Understand the cause
		\item Escalate according to defined procedure
		\item Raising thresholds hides the signal}
	{Validate, understand, escalate.}
	
	\hq{68}{Executive Reports}
	{Executive risk reports should BEST focus on:}
	{\item Technical detail
		\item All raw data
		\item Tool configurations
		\item Trends, business impact and status against appetite}
	{D}
	{\item Executives need decision-relevant information
		\item Trends show direction
		\item Appetite comparison shows what needs action}
	{Executives need business-level insight.}
	
	\hq{69}{Register Updates}
	{The risk register should be updated:}
	{\item After audits only
		\item At project end only
		\item When risk changes occur and at regular intervals
		\item Annually only}
	{C}
	{\item Risk is dynamic
		\item Triggers include changes and incidents
		\item Periodic review catches drift}
	{Update on change and on a schedule.}
	
	\hq{70}{Continuous Monitoring}
	{Continuous monitoring BEST provides:}
	{\item Elimination of risk
		\item Lower cost
		\item Timely awareness of changes in risk
		\item Elimination of audits}
	{C}
	{\item It speeds detection
		\item It supports faster response
		\item It complements, not replaces, audits}
	{Continuous monitoring improves timeliness.}
	
	\hq{71}{KRI Versus KPI}
	{A KRI differs from a KPI because a KRI:}
	{\item Measures achievement of performance goals
		\item Measures cost only
		\item Signals changing risk exposure
		\item Is reported only annually}
	{C}
	{\item KPIs measure performance
		\item KRIs warn of rising exposure
		\item Both are metrics but serve different purposes}
	{KRI warns; KPI measures performance.}
	
	\hq{72}{Too Many KRIs}
	{A dashboard has 150 KRIs and nobody acts on them. What is BEST?}
	{\item Report less often
		\item Reduce to a small set of meaningful, actionable indicators
		\item Remove all indicators
		\item Add more indicators}
	{B}
	{\item Too many metrics hide the signal
		\item Focus on key exposures
		\item Each KRI needs an owner and a response}
	{Fewer, better indicators.}
	
	\hq{73}{All Green}
	{All risk indicators have been green for two years. What should the risk practitioner do?}
	{\item Cut the risk budget
		\item Validate the thresholds and the quality of the indicators
		\item Stop reporting
		\item Celebrate and reduce monitoring}
	{B}
	{\item Constant green may mean thresholds are too loose
		\item Indicators may not reflect real exposure
		\item Validation maintains credibility}
	{Perpetual green is suspicious.}
	
	\hq{74}{Board Reporting}
	{Which is BEST for board risk reporting?}
	{\item Raw KRI data dumps
		\item An annual list of all risks
		\item Detailed technical logs
		\item Concise, trend-based information linked to risk appetite}
	{D}
	{\item Boards need clarity
		\item Trends show direction
		\item Appetite comparison supports oversight}
	{Concise and appetite-linked.}
	
	\hq{75}{Missed Incident}
	{A major incident occurred that no KRI signaled. What should be done?}
	{\item Remove the KRIs
		\item Ignore it as unforeseeable
		\item Review the KRIs and add or adjust indicators
		\item Assign blame to the monitoring team}
	{C}
	{\item KRIs should help anticipate events
		\item Gaps indicate missing or weak indicators
		\item Improve the monitoring framework}
	{Missed signals mean improve the indicators.}
	
	\hq{76}{Register Accuracy}
	{Who should review risk register entries for accuracy?}
	{\item No one
		\item Auditors only
		\item The risk owners, with facilitation by the risk practitioner
		\item Vendors}
	{C}
	{\item Owners know their risks
		\item Practitioners facilitate consistency
		\item Accuracy is an ongoing duty}
	{Owners keep their entries current.}
	
	\hq{77}{Escalation Thresholds}
	{Delayed reporting of significant risks is BEST addressed by:}
	{\item Annual reminders
		\item Defined escalation thresholds and timelines
		\item More forms
		\item Longer meetings}
	{B}
	{\item Clear triggers speed escalation
		\item Timelines create accountability
		\item Reminders alone do not change behavior}
	{Define when and to whom to escalate.}
	
	\hq{78}{KRI Data Source}
	{KRI data is BEST sourced from:}
	{\item Self-reports only
		\item Reliable, preferably automated or independent systems
		\item Informal conversations
		\item Estimates from memory}
	{B}
	{\item Reliable data supports decisions
		\item Automation reduces error
		\item Independence reduces bias}
	{Trust data you can verify.}
	
	\hq{79}{Audit Findings}
	{How should audit findings be used in risk management?}
	{\item As input to update risks and the risk register
		\item They replace risk assessments
		\item They should be ignored
		\item They are only for HR}
	{A}
	{\item Findings reveal control weaknesses
		\item They inform likelihood and impact
		\item They complement assessments}
	{Audit findings are risk inputs.}
	
	\hq{80}{Trend Analysis}
	{What is the main value of trend analysis in risk monitoring?}
	{\item It predicts the exact future
		\item It shows the direction and speed of change in risk
		\item It eliminates risk
		\item It replaces thresholds}
	{B}
	{\item Trends show movement toward or away from appetite
		\item They support proactive response
		\item They do not predict certainty}
	{Trends show where risk is heading.}
	
	% ============================================================
	\section{BIA, Recovery and Third Parties}
	% ============================================================
	
	\hq{81}{BIA Purpose}
	{The PRIMARY purpose of a business impact analysis is to:}
	{\item Test the disaster recovery plan
		\item Select backup vendors
		\item Determine criticality of processes and recovery requirements
		\item Evaluate employee performance}
	{C}
	{\item BIA shows what matters most
		\item It sets recovery time and point objectives
		\item Vendor and testing decisions follow}
	{BIA defines what must be recovered and how fast.}
	
	\hq{82}{RTO}
	{The recovery time objective (RTO) is the:}
	{\item Target time to restore a process after disruption
		\item Cost of recovery
		\item Maximum acceptable data loss
		\item Maximum tolerable downtime}
	{A}
	{\item RTO is a recovery time target
		\item RPO is about data loss
		\item MTD is the absolute limit}
	{RTO is about time to recover.}
	
	\hq{83}{RPO}
	{The recovery point objective (RPO) describes:}
	{\item Total downtime
		\item The maximum acceptable data loss measured in time
		\item Cost of recovery
		\item Time to restore service}
	{B}
	{\item RPO sets how far back data can be restored
		\item It drives backup frequency
		\item RTO is separate}
	{RPO is about data loss.}
	
	\hq{84}{RTO and MTD}
	{Which relationship is correct?}
	{\item RTO should be longer than MTD
		\item RTO and MTD must always be equal
		\item They are unrelated
		\item RTO should be shorter than the maximum tolerable downtime}
	{D}
	{\item MTD is the point beyond which impact becomes unacceptable
		\item RTO is a target inside that limit
		\item RTO beyond MTD defeats the purpose}
	{RTO sits inside MTD.}
	
	\hq{85}{BIA Input}
	{Recovery priorities in a BIA should be determined primarily by:}
	{\item Internal audit
		\item IT staff only
		\item Business process owners
		\item Vendors}
	{C}
	{\item Business owners understand impact
		\item IT validates feasibility
		\item Priorities reflect business needs}
	{The business sets priorities; IT delivers.}
	
	\hq{86}{Backup Gap}
	{The RPO is 4 hours but backups run nightly. What is the issue?}
	{\item The RTO is too short
		\item Backup frequency cannot meet the RPO
		\item Backups run too often
		\item There is no issue}
	{B}
	{\item Nightly backups can lose up to 24 hours of data
		\item The RPO requires more frequent recovery points
		\item Align backup design to RPO}
	{Backups must match the RPO.}
	
	\hq{87}{Failed DR Test}
	{A disaster recovery test restores a system in 10 hours against an RTO of 4 hours. What should happen?}
	{\item Ignore the result
		\item Cancel future tests
		\item Raise the RTO without review
		\item Investigate causes, improve the plan and retest}
	{D}
	{\item The plan does not meet the objective
		\item Root causes must be addressed
		\item Retesting confirms improvement
		\item Changing the RTO needs business agreement}
	{A failed test is a finding, not a formality.}
	
	\hq{88}{Vendor Assessment Timing}
	{When should a vendor\textquotesingle s risk be assessed?}
	{\item Only after go-live
		\item After an incident
		\item Never, if the vendor is reputable
		\item Before contract signature and periodically after}
	{D}
	{\item Pre-contract assessment supports negotiation
		\item Risk changes over time
		\item Reputation is not an assessment}
	{Assess before you sign.}
	
	\hq{89}{Right to Audit}
	{Which contract provision BEST supports third-party risk oversight?}
	{\item Right-to-audit and security reporting clauses
		\item A price discount clause
		\item A logo usage clause
		\item A marketing clause}
	{A}
	{\item Oversight needs access to evidence
		\item Reporting duties provide ongoing visibility
		\item Commercial terms do not provide assurance}
	{Contracts must enable oversight.}
	
	\hq{90}{Vendor Concentration}
	{Many critical processes depend on one vendor. What is the BEST response?}
	{\item Increase reliance further
		\item Stop monitoring the vendor
		\item Assess concentration risk and plan alternatives or contingency
		\item Ignore it}
	{C}
	{\item Single dependency creates concentrated exposure
		\item Alternatives or exit plans reduce impact
		\item Monitoring should increase}
	{Concentration is a risk in itself.}
	
	% ============================================================
	\section{Governance and Risk Culture}
	% ============================================================
	
	\hq{91}{Register Purpose}
	{The main purpose of the risk register is to:}
	{\item Record and track risks, owners and responses
		\item Replace audits
		\item Store incident tickets
		\item Track project budgets}
	{A}
	{\item It provides a central view of risks
		\item Owners and responses support accountability
		\item It supports reporting}
	{The register supports tracking and accountability.}
	
	\hq{92}{Risk Owner Qualities}
	{A risk owner should have:}
	{\item Technical skills only
		\item Audit independence
		\item Accountability and authority to manage the risk
		\item No decision authority}
	{C}
	{\item Owners make response decisions
		\item Authority matches accountability
		\item Technical skills can be supplied by others}
	{Ownership requires authority.}
	
	\hq{93}{Risk Culture}
	{Risk culture is MOST influenced by:}
	{\item The tools used
		\item The number of policies
		\item Annual training alone
		\item The tone and behavior of senior leadership}
	{D}
	{\item People follow what leaders do
		\item Consistent behavior signals priorities
		\item Policies and tools do not create culture alone}
	{Tone at the top shapes culture.}
	
	\hq{94}{Three Lines}
	{In the three lines model, who owns and manages risk day to day?}
	{\item Regulators
		\item Internal audit
		\item The risk function
		\item Business and operational management (first line)}
	{D}
	{\item First line owns and manages risk
		\item Second line oversees and advises
		\item Third line provides independent assurance}
	{The first line owns the risk.}
	
	\hq{95}{Practitioner Role}
	{The risk practitioner\textquotesingle s role in risk acceptance is to:}
	{\item Advise and ensure an informed decision by the risk owner
		\item Accept risk on behalf of the business
		\item Hide residual risk
		\item Reject all risk acceptance}
	{A}
	{\item Practitioners provide analysis and advice
		\item Owners decide
		\item Decisions should be informed and documented}
	{Advise; do not decide.}
	
	\hq{96}{Emerging Technology}
	{Before adopting an emerging technology, the organization should:}
	{\item Adopt it widely at once
		\item Perform a risk assessment
		\item Wait for incidents
		\item Always avoid new technology}
	{B}
	{\item New technology brings new risks
		\item Assessment informs the decision
		\item Blanket avoidance can hurt the business}
	{Assess before adopting.}
	
	\hq{97}{Assessment Frequency}
	{Risk assessment frequency should BEST be:}
	{\item Periodic and triggered by significant change
		\item After incidents only
		\item Annual only
		\item During audits only}
	{A}
	{\item Periodic reviews catch drift
		\item Changes and incidents trigger reassessment
		\item Annual is a minimum, not a maximum}
	{Periodic plus event-driven.}
	
	\hq{98}{Communicating Risk}
	{Risk should be communicated to non-technical executives in terms of:}
	{\item Business impact such as financial, operational and reputational effects
		\item Technical jargon
		\item Vulnerability identifiers
		\item Protocol details}
	{A}
	{\item Executives decide on business terms
		\item Impact links to objectives
		\item Technical detail can obscure the message}
	{Speak the language of the business.}
	
	\hq{99}{Stale Register}
	{Many register entries have not been reviewed in over two years. What is BEST?}
	{\item Re-validate entries with owners and establish a review cadence
		\item Leave them as they are
		\item Increase all ratings
		\item Delete all entries}
	{A}
	{\item Stale entries may be wrong
		\item Owners must confirm status
		\item A cadence prevents recurrence}
	{Stale registers mislead decisions.}
	
	\hq{100}{Unclear Ownership}
	{A risk has no clear owner. What should the risk practitioner do?}
	{\item Escalate so that an accountable owner is assigned
		\item Leave it blank
		\item Assume IT owns it
		\item Close the risk}
	{A}
	{\item Unowned risks are unmanaged
		\item Ownership must be assigned by appropriate authority
		\item Defaulting to IT confuses accountability}
	{Every risk needs an owner.}
	
	\begin{frame}{Exam Mindset Summary}
		\begin{block}{Remember}
			\begin{itemize}
				\item \textbf{FIRST:} understand context and assess before acting
				\item \textbf{Owners decide:} the business owner accepts risk; the practitioner advises
				\item \textbf{Appetite:} compare residual risk to appetite, then escalate or treat
				\item \textbf{Accountability} cannot be outsourced or transferred
				\item \textbf{Monitor:} validate data, trend it, and keep the register current
			\end{itemize}
		\end{block}
	\end{frame}
	
	\begin{frame}{Thank You}
		\begin{center}
			{\Huge \textbf{Good Luck on Your CRISC Exam!}}
		\end{center}
	\end{frame}
	
\end{document}